\documentclass[11pt]{letter}

\usepackage[margin=1in]{geometry}
\usepackage{hyperref}

\signature{Oliver Pardo}
\address{Departamento de Administraci\'on\\
Pontificia Universidad Javeriana\\
Bogot\'a, Colombia\\
\href{mailto:pardoo@javeriana.edu.co}{pardoo@javeriana.edu.co}}

\begin{document}

\begin{letter}{Editors\\
Journal of Public Economic Theory}

\opening{Dear Editors,}

Please consider my manuscript, ``On the efficiency of mandatory retirement savings under endogenous borrowing constraints,'' for publication in the \textit{Journal of Public Economic Theory}.

The paper studies mandatory retirement saving in a life-cycle model in which borrowing constraints arise endogenously from limited pledgeability and moral hazard. The main result is that mandatory contributions can reduce welfare when they lock resources away from young borrowers who need credit to undertake positive-net-present-value investments. The paper also shows that this distortion can be mitigated when retirement balances are pledgeable, and that a pension design with unconditional basic savings can relax borrowing constraints for low-wealth individuals.

I believe the manuscript is a good fit for the \textit{Journal of Public Economic Theory} because it links pension design, collateral constraints, and credit-market inefficiency within a simple public-economic-theory framework. The paper contributes to the theoretical analysis of retirement saving policy by emphasizing that the welfare effect of mandatory contributions depends not only on forced saving, but also on whether accumulated balances can be used as collateral.

The manuscript is original, has not been published, and is not under consideration elsewhere. I have no conflicts of interest to disclose.

Thank you for considering the manuscript. I look forward to hearing from you.

\closing{Sincerely,}

\end{letter}

\end{document}
